\documentclass[12pt]{article}

%Importamos paquetes
\usepackage[utf8]{inputenc}  %Configura caracteres especiales
\usepackage[spanish]{babel}  %Configura en español los valores predetermina
\usepackage{amsmath}   %Para el manejo de formulas
\usepackage{array}  %Mejor manejo de tablas
\usepackage{geometry} %Margenes

\newgeometry{margin=2.5cm} %Modifica el margen de la hoja

\newcommand{\minombre}{Jesús Aguilar}
\newcommand{\nombrepractica}{Práctica 4: tablas , entornos y comandos}
\newcommand{\universidad}{UASLP-FEPZM}




\newenvironment{ejercicio}
{
   \vspace{0.3cm}
   \noindent\textbf{Ejercicio:}
   \begin{itshape}  %Modifica a italica el texto

}
{
\end{itshape}
\vspace{0.3cm}
}
\begin{document}
\begin{center}
\vspace*{4cm}
{\LARGE \nombrepractica \LaTeX}\\[0.5cm]
\vspace{2.5cm}
{\large \minombre}\\[0.5cm]
\vspace{2.5cm}
{\large \today}
\end{center}

\section{Intrdducción}
En esta practica aprenderemos a crear tablas y a usar \textbf{enviroments}, inclueyndo enviroments personalizados

\section{Uso de tablas}
\begin{table}[h]
    \centering
    \begin{tabular}{c|c|c}
         %\hline
         Alumno&Materia&Calificacion\\
         \hline
         Pedro&MTeologia&8\\
         Delfina&Biologia&7
    \end{tabular}{}
    \caption{Calificaciones de alumnos}
    \label{tab:calificaciones}

\end{table}

\subsection{Tablas con columnas combinadas}

\begin{table}[h]
    \centering
    \begin{tabular}{c|c|c}
         \hline
         \multicolumn{3}{|c|}{Horario}\\ \hline
         Dia&Hora&Materia \\ \hline
         Lunes&10-11&Matematicas \\ \hline
         Martes&11-12&fISICA \\ \hline
    \end{tabular}
    \caption{Horario semanal}
    \label{tab:Horario}
\end{table}
\section{Enviroment perosnalizado}
Ahora usaremos nuestro enviroment \textbf{ejercicio}

\begin{ejercicio}
    Crear una tabla con al menos 4
    
\end{ejercicio}


\end{document}

